\section{Test Case 21: Anchored Dynamic Refinement with Platen Evaporation}

Test Case 21 validates the hybrid OpenACC port of dynamic Akima refinement.
It combines Maxwell rheology, the stochastic Platen integrator,
Yarin evaporation, aerodynamic drag, nozzle insertion, and numerical anchor
beads. The input is stored in
\texttt{examples/input-21/input.dat}.

The initial jet contains 400 elements at a resolution of 0.02 cm and occupies
the interval from the nozzle to 8 cm, while the collector remains at 16 cm.
The free half of the domain permits field-driven stretching. Interior beads
are tagged as anchors every 0.10 cm. These beads are ordinary material points
that remain knots of the old and refined interpolation meshes; no variable
mass or nanoparticle model is enabled.

The refinement target is also 0.10 cm. The reduced charge density limits
sensitivity to the order of the direct Coulomb sum, whereas the increased
external potential makes the short validation reach the refinement threshold.
These choices accelerate a numerical validation and are not proposed as a
reference experimental parameter set. The minimum interval between accepted
remeshes is $5\times10^{-5}$ s, leaving a post-refinement trajectory before
the run ends.

The OpenACC path keeps the Platen stages, pre-generated Gaussian history,
insertion, and state on the device. Repeated refinement checks transfer only
three reduction scalars. At an accepted event the active state is downloaded,
the host constructs normalized source and target knots, and Akima coefficient
construction and all 11 spline interpolations execute on the GPU. Conservation
and anchor restoration remain on the host, after which the remeshed state is
uploaded once. The
initial allocation reserves 100 additional entries, so the validated event
does not reallocate a mapped array. An additional capacity-growth mode reduces
only this developer reserve to 50 entries. It releases the old persistent
mappings before the host allocation changes, expands the Gaussian history and
Platen/Coulomb workspaces, and binds the completed mesh once. The same device
Akima kernels are used in both modes.

With GFortran the event occurs at step 7,134 and changes the active count from
413 to 454, ending with 455 elements. With NVFORTRAN 24.3 the CPU event occurs
at step 7,118 and changes 412 elements to 460, ending with 462. The complete
host-force oracle on one NVIDIA A30 reproduces the NVFORTRAN CPU event and
topology. The native A30 path accepts at step 7,122, changes 412 elements to
459, and ends with 461. The four-step shift is consistent with the different
target-centric Coulomb summation order in a bending-sensitive trajectory.

All paths start with 79 interior anchors, finish with \texttt{nref=1}, and
preserve anchor positions and fields within the $10^{-12}$ checker tolerance.
Reference and evaporation-volume conservation errors are of order
$10^{-16}$, with no capacity reallocation. The 8,000-step run pre-generates
24,048,000 Gaussian doubles for the full 500-entry capacity and uploads them
once. A transfer audit observed only the three reduction scalars during
ordinary threshold checks, one full-state download and upload at the Akima
event, and the independent final-shutdown download.

In the capacity-growth mode, NVFORTRAN CPU accepts at step 7,156 and the
native A30 at step 7,155. Both change 413 active elements to 455 and finish
with 456. The complete-force oracle accepts at the CPU step with the same
topology. Capacity grows from 450 to 557 entries. Existing Gaussian values are
preserved while the history is repacked, new indices are generated once, and
26,784,000 doubles are uploaded for the new stride. CPU and oracle statistics
agree within $3\times10^{-7}$ relatively; the native GPU statistics differ by
at most 0.8 percent while retaining identical topology and conservation
invariants. A transfer audit again found only the accepted-event and final
full-state downloads, plus one topology/evaporation rebind, one history
upload, and one workspace recreation at the event.

The script \texttt{tests/refinement/run.sh} defaults to GFortran. Its
\texttt{nvfortran}, \texttt{openacc}, and \texttt{force-oracle} arguments run
the corresponding event-level validation paths. Appending
\texttt{capacity-growth} enables the forced-capacity validation.
